% =====================================================================================
% Subsection "Quality of the Certified Model", \input at the end of evaluation.tex. Condensed from
% generation_quality.tex (Section 6.3); sources docs/improvements/F2_quality.md, the scripts
% experiments/2_attack/real_weights_ppl.py and 9_quality/ (window setting). The attribution table,
% smoothing details and the full literature comparison are in ext_generation.tex.
% Labels defined: sec:quality, tab:quality (bottom half of the draft's table only).
% Labels used: sec:intmodel (defence.tex).
% =====================================================================================

\providecommand{\para}[1]{\par\noindent\textbf{#1}\quad}
\providecommand{\todo}[1]{\textbf{[TODO: #1]}}

\subsection{Quality of the Certified Model}
\label{sec:quality}

The provider must serve the certified integer model (Section~\ref{sec:intmodel}), with one scale
per weight matrix and static activation scales. We quantise the real OPT
checkpoints~\cite{zhang2022opt} and measure perplexity on the first 16 non-overlapping 128-token
windows of the WikiText-2 (raw) test split~\cite{merity2017wikitext}, each scored on its own
(2{,}032 predicted tokens), calibrating on the split's last 128-token window. Published results
use 2{,}048-token sequences, so we compare ratios to fp32 on the same windows.

\para{Where the perplexity goes.} The integer LayerNorm outputs a gain times the normalised value.
With the best scalar gain, perplexity is 17\% above fp32 for OPT-125M, 11\% for OPT-1.3B
(Table~\ref{tab:quality}) and, in a separate run, 33\% for OPT-6.7B (36.0 against 27.1).
Switching each component of a floating-point copy of the graph between integer and float locates
OPT-125M's loss: float LayerNorm outputs remove more than half of the gap (75.5 to 69.4, against
64.8 in fp32), while the 8-bit softmax, the residual stream and the ReLU table each cost at most
0.1 (extended version~\cite{fullversion}). The loss comes from rounding, not clipping: a few
outlier channels of OPT's residual stream inflate the standard deviation, so a scalar gain that
keeps them within int8 leaves the other channels few levels.

\begin{table}[t]
  \caption{Integer OPT models on WikiText-2 (16 windows of 128 tokens): perplexity (ratio to fp32)
  with the best scalar norm gain and with smoothed gains ($\alpha=0.55$), and top-1 agreement with
  fp32 before and after smoothing.}
  \label{tab:quality}
  \centering\footnotesize
  \setlength{\tabcolsep}{4pt}
  \begin{tabular}{@{}lrrrc@{}}
    \toprule
    Model & fp32 & Scalar gain & Smoothed & Top-1 agr. \\
    \midrule
    OPT-125M & 64.76 & 75.52 ($\times$1.17) & 66.62 ($\times$1.03) & 0.72$\to$0.82 \\
    OPT-1.3B & 32.79 & 36.44 ($\times$1.11) & 33.70 ($\times$1.03) & 0.82$\to$0.90 \\
    \bottomrule
  \end{tabular}
\end{table}

\para{SmoothQuant through the norm's gains.} Per-channel gains absorb SmoothQuant's
scaling~\cite{xiao2023smoothquant}: channel $j$ gets the gain $\gamma_j=\kappa/s_j$ with
$s_j=\rho_j^{\alpha}/\omega_j^{1-\alpha}$, where $\rho_j$ is the 99.9th percentile of the channel's
normalised magnitude on the calibration window and $\omega_j$ the largest magnitude in column $j$
of the matrices that the norm feeds, into which $1/\gamma_j$ is folded. This is SmoothQuant-O3
(per-tensor weights, static per-tensor activations) inside the integer graph. Multipliers stay
scalars and no operation or shape changes, so the costs, the verifier and V1 are unaffected; the
graph digest binds the gains. With $\alpha=0.55$ and 99.99th-percentile clipping at the q/k/v,
fc1 and attention-output requantisations, both models reach $\times1.03$; over 32 windows and
three calibration windows, OPT-125M stays within $\times1.03$--$1.05$. Per-row power-of-two weight
scales reach $\times1.00$--$1.01$ on OPT-125M but need four multipliers per requantisation, which
V1 does not support yet.

\para{Comparison.} SmoothQuant-O3 raises perplexity by 1.6\% on OPT-175B and 0.8\% on
OPT-IML-30B~\cite{xiao2023smoothquant}; ZeroQuant-V2, with per-token activations and 16-bit
weights, by about 1\% on OPT-125M and 4\% on OPT-1.3B~\cite{yao2023zeroquantv2}. Proof systems
that keep 16-bit values stay under 1\%~\cite{zkllm,zkgpt}. We set targets of $\times1.06$,
$\times1.04$ and $\times1.05$ for OPT-125M, OPT-1.3B and OPT-6.7B; the two smaller models meet
theirs. \todo{smoothed OPT-6.7B, Qwen3-4B} \todo{full test set at 2{,}048 tokens, incl.\ 8-bit
attention probabilities} \todo{real-weight run of the smoothed graphs, incl.\ V1}
